\section{Introduction}
\label{sec:intro}

Swarms of unmanned aerial vehicles (UAVs) are used for area surveillance, disaster response and data
collection, where no ground infrastructure is available. The UAVs then relay each other's traffic over a
flying ad hoc network (FANET)~\cite{bekmezci2013fanet}. FANET nodes move fast and in three dimensions, so
the radio link between two neighbours often lasts only a few seconds. Most routing protocols notice a
broken link only after the fact, when a transmission fails or a neighbour timer expires. By then the packet
in transit has already been lost or delayed, and the route must be rebuilt~\cite{khan2020survey}.

Greedy forwarding~\cite{karp2000gpsr} sends each packet to the neighbour closest to the destination. That
neighbour is usually near the edge of radio range, so its link is the one most likely to break next. On UavNetSim-v1~\cite{zhou2025uavnetsim,zhou2026repo}, the
platform used in this work, the medium access control (MAC) layer then retries the same neighbour up to
five times before dropping the packet, and a neighbour stays in the table for up to 2\,s after its
last hello message.

This paper asks two questions. First, how many packets do the existing FANET routing protocols actually
lose because a link breaks, rather than because of a MAC collision or another cause? Second, can a node
that predicts the remaining lifetime of each link avoid those losses, and how much of the gain comes from
the prediction itself rather than from reacting faster?

Link-lifetime prediction has been studied before. The link expiration time (LET) of Su~\emph{et~al.}~\cite{su2001mobility}
predicts when two nodes moving at constant velocity leave each other's range, and OPAR~\cite{gharib2021opar}
uses it in a centralised route computation. Learning-based protocols such as QGeo~\cite{jung2017qgeo},
Q-FANET~\cite{dacosta2021qfanet} and TARRAQ~\cite{cui2022tarraq} adapt routes from observed delivery outcomes.
What is usually missing is a measurement of where packets are lost, and a test that separates the effect of
prediction from the effect of any reaction to a failure.

Our contributions are:
\begin{itemize}
  \item An exact loss attribution for UavNetSim-v1. Because the platform's physical layer is deterministic,
  each failed reception can be classified as out of range (link break) or interference (collision). We apply
  it to seven protocols over 231 runs.
  \item LEAP-GR, a distributed geographic protocol that behaves exactly like Greedy forwarding unless the
  chosen neighbour is predicted to leave range within a threshold $\tauth$, in which case it switches to a
  neighbour predicted to stay.
  \item A paired evaluation with an identical MAC that compares LEAP-GR with Greedy, GRAd and a reactive
  control that switches next hop only after a failure, so that the gain of prediction alone is isolated.
\end{itemize}

\Cref{sec:related} reviews related work. \Cref{sec:model} presents the system model, the loss attribution
and LEAP-GR. \Cref{sec:eval} describes the experimental setup and analyses the results.
\Cref{sec:conclusion} concludes and lists the limits of this study.
